% from ega2-1-fr.tex
\oldpage[II]{15}
à un isomorphisme canonique près~; en outre, $\mathbf{S}(E)$ est un foncteur
covariant en $E$, de la catégorie des $A$-modules dans celle des $A$-algèbres
commutatives~; enfin, la caractérisation précédente montre aussi que si
$E=\varinjlim E_\lambda$, on a
$\mathbf{S}(E)=\varinjlim\mathbf{S}(E_\lambda)$ à un isomorphisme canonique
près. Par abus de langage, un produit
$\sigma(x_1)\sigma(x_2)\cdots\sigma(x_n)$, où les $x_i\in E$, se note souvent
$x_1x_2\cdots x_n$ si aucune confusion n'en résulte. L'algèbre
$\mathbf{S}(E)$ est \emph{graduée}, $\mathbf{S}_n(E)$ étant l'ensemble des
combinaisons linéaires des produits de $n$ éléments de $E$ ($n\geq 0$)~;
l'algèbre $\mathbf{S}(A)$ est canoniquement isomorphe à l'algèbre des
polynômes $A[T]$ à une indéterminée et l'algèbre $\mathbf{S}(A^n)$ à
l'algèbre des polynômes à $n$ indéterminées $A[T_1,\ldots,T_n]$.
\end{env}

\begin{env}[1.7.2]
\label{II.1.7.2-fr}
Soit $\varphi$ un homomorphisme d'anneaux $A\to B$. Si $F$ est un $B$-module,
l'application canonique $F\to\mathbf{S}(F)$ donne une application canonique
$F_{[\varphi]}\to\mathbf{S}(F)_{[\varphi]}$, qui se factorise donc en
$F_{[\varphi]}\to\mathbf{S}(F_{[\varphi]})\to
\mathbf{S}(F)_{[\varphi]}$~; l'homomorphisme canonique
$\mathbf{S}(F_{[\varphi]})\to\mathbf{S}(F)_{[\varphi]}$ est surjectif, mais
non nécessairement bijectif. Si $E$ est un $A$-module, tout di-homomorphisme
$E\to F$ (c'est-à-dire tout $A$-homomorphisme $E\to F_{[\varphi]}$) donne donc
canoniquement un $A$-homomorphisme d'algèbres
$\mathbf{S}(E)\to\mathbf{S}(F_{[\varphi]})\to
\mathbf{S}(F)_{[\varphi]}$, c'est-à-dire un di-homomorphisme d'algèbres
$\mathbf{S}(E)\to\mathbf{S}(F)$.

Avec les mêmes notations, pour tout $A$-module $E$,
$\mathbf{S}(E\otimes_A B)$ s'identifie canoniquement à l'algèbre
$\mathbf{S}(E)\otimes_A B$~; cela résulte aussitôt de la propriété universelle
de $\mathbf{S}(E)$ (\hyperref[II.1.7.1-fr]{1.7.1}).
\end{env}

\begin{env}[1.7.3]
\label{II.1.7.3-fr}
Soit $R$ une partie multiplicative de l'anneau $A$~; appliquant
(\hyperref[II.1.7.2-fr]{1.7.2}) à l'anneau $B=R^{-1}A$, et se rappelant que
$R^{-1}E=E\otimes_A R^{-1}A$, on voit que l'on a
$\mathbf{S}(R^{-1}E)=R^{-1}\mathbf{S}(E)$ à un isomorphisme canonique près.
En outre, si $R'\supset R$ est une seconde partie multiplicative de $A$, le
diagramme
\[
  \xymatrix{
    R^{-1}E\ar[r]\ar[d] & {R'}^{-1}E\ar[d]\\
    \mathbf{S}(R^{-1}E)\ar[r] & \mathbf{S}({R'}^{-1}E)
  }
\]
est commutatif.
\end{env}

\begin{env}[1.7.4]
\label{II.1.7.4-fr}
Soit maintenant $(S,\mathcal{A})$ un espace annelé, et soit $\mathcal{E}$ un
$\mathcal{A}$-Module sur $S$. Si, à tout ouvert $U\subset S$ on associe le
$\Gamma(U,\mathcal{A})$-module
$\mathbf{S}(\Gamma(U,\mathcal{E}))$, on définit (vu le caractère fonctoriel
de $\mathbf{S}(E)$ (\hyperref[II.1.7.2-fr]{1.7.2})) un préfaisceau d'algèbres~;
on dit que le faisceau associé, que l'on note $\mathbf{S}(\mathcal{E})$ ou
$\mathbf{S}_{\mathcal{A}}(\mathcal{E})$ est la $\mathcal{A}$-Algèbre
\emph{symétrique} du $\mathcal{A}$-Module $\mathcal{E}$. Il résulte aussitôt
de (\hyperref[II.1.7.1-fr]{1.7.1}) que $\mathbf{S}(\mathcal{E})$ est solution
d'un problème universel~: tout homomorphisme de $\mathcal{A}$-Modules
$\mathcal{E}\to\mathcal{B}$, où $\mathcal{B}$ est une $\mathcal{A}$-Algèbre,
se factorise d'une seule manière en
$\mathcal{E}\to\mathbf{S}(\mathcal{E})\to\mathcal{B}$, la seconde flèche étant
un homomorphisme de $\mathcal{A}$-Algèbres. Il y a donc correspondance
biunivoque entre homomorphismes $\mathcal{E}\to\mathcal{B}$ de
$\mathcal{A}$-Modules et homomorphismes
$\mathbf{S}(\mathcal{E})\to\mathcal{B}$ de $\mathcal{A}$-Algèbres. En
particulier, tout homomorphisme $u:\mathcal{E}\to\mathcal{F}$ de
$\mathcal{A}$-Modules définit un homomorphisme
$\mathbf{S}(u):\mathbf{S}(\mathcal{E})\to\mathbf{S}(\mathcal{F})$ de
$\mathcal{A}$-Algèbres et $\mathbf{S}(\mathcal{E})$ est donc un foncteur
covariant en $\mathcal{E}$.

